%! The Secant, Cosecant, and Cotangent Functions — Unit 7, Lesson 7.4 — Lesson Plan
\documentclass[10pt]{article}
\usepackage{precalculus-article}
\usepackage{precalculus-boxes}
\usepackage{graphicx}
\graphicspath{{images/}}
\newcommand{\TallMath}[1]{$\displaystyle #1\rule[-1.4em]{0pt}{3.2em}$}
\newcommand{\CourseName}{Precalculus}
\newcommand{\MeetingLength}{55 minutes}
\newcommand{\UnitNumberName}{Unit 7: Analytic Trigonometry and Polar Coordinates \quad}
\newcommand{\LessonNumberName}{Lesson 7.4: The Secant, Cosecant, and Cotangent Functions}

\begin{document}

%----------------------------------------------------------------------
% Title Block
%----------------------------------------------------------------------
\begin{center}
  {\Large\bfseries \CourseName} \\[4pt]
  {\small\bfseries \UnitNumberName \LessonNumberName \quad (2 instructional periods, \MeetingLength\ each)}
\end{center}

%----------------------------------------------------------------------
% Primary Objective
%----------------------------------------------------------------------
\begin{tcolorbox}[colback=sky, colframe=navy, boxrule=0.9pt, arc=2mm,
    left=2mm, right=2mm, top=1.5mm, bottom=1.5mm]
  \textbf{Primary Objective:} Students will identify the key characteristics
  of the secant, cosecant, and cotangent functions --- including their definitions
  as reciprocals, their domains, their vertical asymptotes, and their ranges ---
  by connecting each reciprocal function to its parent sine, cosine, or tangent
  function.
  \textit{(Big Idea: Functions)}
\end{tcolorbox}

\vspace{0.08in}

%----------------------------------------------------------------------
% Priority Ideas & Skills
%----------------------------------------------------------------------
\begin{skillbox}[Priority Ideas \& Skills]{goldbox}
\begin{tabularx}{\linewidth}{@{} p{0.46\linewidth} X @{}}
\textbf{AP Skills} & \textbf{Key Understandings (EKs)} \\
\midrule
\textbf{2.B} \textit{Construct equivalent representations} ---
  Express sec, csc, and cot as quotients of sine and cosine, and identify
  which parent-function zero causes each asymptote.
&
\textbf{EK 3.11.A.1}: $\sec\theta = \dfrac{1}{\cos\theta}$, defined where
$\cos\theta \neq 0$. \\[6pt]

\textbf{3.A} \textit{Describe characteristics} ---
  State the domain, range, period, and vertical asymptote locations for
  each of the three reciprocal trigonometric functions.
&
\textbf{EK 3.11.A.2}: $\csc\theta = \dfrac{1}{\sin\theta}$, defined where
$\sin\theta \neq 0$. \\[6pt]

&
\textbf{EK 3.11.A.3}: The graphs of sec and csc have vertical asymptotes
where $\cos\theta = 0$ and $\sin\theta = 0$, respectively; range
$= (-\infty,-1]\cup[1,\infty)$. \\[6pt]

&
\textbf{EK 3.11.A.4}: $\cot\theta = \dfrac{\cos\theta}{\sin\theta}$
(equivalently, $\dfrac{1}{\tan\theta}$ where $\tan\theta \neq 0$),
defined where $\sin\theta \neq 0$. \\[6pt]

&
\textbf{EK 3.11.A.5}: The graph of cot has vertical asymptotes where
$\sin\theta = 0$ and is decreasing on each interval between consecutive
asymptotes. \\
\end{tabularx}
\end{skillbox}

\vspace{0.08in}

%----------------------------------------------------------------------
% Vocabulary
%----------------------------------------------------------------------
\begin{skillbox}[{Vocabulary, Concepts, \& Theorems}]{greenbox}
\begin{tabularx}{\linewidth}{@{} p{0.24\linewidth} X @{}}
\textbf{Term} & \textbf{Definition / Note} \\
\midrule
secant ($\sec\theta$) &
  $\sec\theta = 1/\cos\theta$; the multiplicative reciprocal of cosine.
  Undefined when $\cos\theta = 0$, i.e., when $\theta = \pi/2 + n\pi$. \\[4pt]
cosecant ($\csc\theta$) &
  $\csc\theta = 1/\sin\theta$; the multiplicative reciprocal of sine.
  Undefined when $\sin\theta = 0$, i.e., when $\theta = n\pi$. \\[4pt]
cotangent ($\cot\theta$) &
  $\cot\theta = \cos\theta/\sin\theta$; the reciprocal of tangent
  (where $\tan\theta \neq 0$).
  Undefined when $\sin\theta = 0$. \\[4pt]
vertical asymptote &
  A vertical line $\theta = a$ where the function value grows without
  bound; caused by a zero in the denominator. \\[4pt]
reciprocal function &
  A function formed by taking $1/f(\theta)$; zeros of $f$ become
  asymptotes of the reciprocal, and $f(\theta) = \pm 1$ are invariant
  points. \\
\end{tabularx}
\end{skillbox}

\vspace{0.08in}

%----------------------------------------------------------------------
% Activate Prior Knowledge
%----------------------------------------------------------------------
\begin{skillbox}[Activate Prior Knowledge \& Spiral Review]{sky}
\begin{tabularx}{\linewidth}{@{}X c@{}}
\begin{minipage}[t]{0.96\linewidth}\vspace{0pt}
\textbf{Spiral topics activated during Warm-Up (first 8--10 minutes):}
\begin{itemize}[leftmargin=1.4em, itemsep=2pt]
  \item \textbf{Unit-circle values for sin and cos}: Students evaluate
    $\sin\theta$ and $\cos\theta$ at $0,\,\pi/6,\,\pi/4,\,\pi/3,\,\pi/2$
    and their quadrant-II/III/IV equivalents --- essential for computing
    sec, csc, and cot by hand.
  \item \textbf{Tangent as a quotient}: $\tan\theta = \sin\theta/\cos\theta$.
    Students must recall this ratio formula; cotangent flips it.
  \item \textbf{Asymptote behavior of $\tan\theta$}: Asymptotes at
    $\theta = \pi/2 + n\pi$ (where $\cos\theta = 0$) foreshadow the
    asymptote structure of secant.
  \item \textbf{Reciprocal arithmetic}: For any nonzero $x$, $1/x$ and $x$
    share the sign, and $|1/x| \geq 1$ when $|x| \leq 1$.
\end{itemize}
\end{minipage}
&
\end{tabularx}
\end{skillbox}

\vspace{0.08in}

%----------------------------------------------------------------------
% Hook
%----------------------------------------------------------------------
\begin{skillbox}[Hook --- Entry Event]{sky}
\textbf{Scenario:} Project (or sketch on the board) a portion of the graph
of $y = \sec\theta$ over $[0, 2\pi]$. The graph shows two ``U-shaped'' arcs
with a gap at $\theta = \pi/2$ and $\theta = 3\pi/2$, where vertical
dashed lines mark the asymptotes.

\textbf{Discussion prompts:}
\begin{enumerate}[leftmargin=1.4em, itemsep=3pt]
  \item \textit{``This is a trigonometric function. What familiar function does
    it look related to? What causes the gaps?''}
  \item \textit{``The output values 0.5 and $-0.5$ never appear. Why might
    the output skip values between $-1$ and $1$?''}
  \item \textit{``Can a trig function have a range that is NOT an interval?
    Have you seen that before?''}
\end{enumerate}

\textbf{Key tension to surface:} We have studied sine, cosine, and tangent.
Each has a ``partner'' function built by taking reciprocals. The zeros of
the parent become the asymptotes of the reciprocal, and the invariant points
($y = \pm 1$) act as the ``bridges'' between the two families.

\textbf{Transition:} ``Today we'll define secant, cosecant, and cotangent
as reciprocal functions, find their domains and asymptotes, and describe
the key features of each graph.''
\end{skillbox}

\vspace{0.08in}

%----------------------------------------------------------------------
% Lesson
%----------------------------------------------------------------------
\begin{skillbox}[Lesson --- Instruction Sequence]{sky}
\begin{multicols}{2}

\textbf{Step 1 --- Secant from Cosine} (12 min)

Define $\sec\theta = 1/\cos\theta$. At each unit-circle angle, if
$\cos\theta = 0$, sec is undefined; if $\cos\theta = 0.5$, $\sec\theta = 2$.
Build a small table of values (use guided notes table) and note that
$\sec\theta \geq 1$ or $\sec\theta \leq -1$ always.
Point out the ``U-arcs'' that open away from the $\theta$-axis, with vertices
at the maximum/minimum of cosine. Formalize EK 3.11.A.1 and 3.11.A.3.
Key question: \textit{``Where are the asymptotes? What equation describes them?''}
($\theta = \pi/2 + n\pi$.)

\columnbreak

\textbf{Step 2 --- Cosecant from Sine} (10 min)

Parallel development: define $\csc\theta = 1/\sin\theta$. Build the table
together; asymptotes where $\sin\theta = 0$ ($\theta = n\pi$). Same range
$(-\infty,-1]\cup[1,\infty)$. Formalize EK 3.11.A.2 and 3.11.A.3.
Emphasize the structural symmetry: the two functions are horizontal
translates of each other by $\pi/2$.

\end{multicols}

\begin{multicols}{2}

\textbf{Step 3 --- Cotangent from Tangent} (10 min)

Define $\cot\theta = \cos\theta/\sin\theta$. Unlike sec and csc, cotangent
is NOT bounded away from zero; its range is all reals.
Asymptotes where $\sin\theta = 0$ (same as csc).
On each interval between asymptotes, cot is \emph{decreasing} (contrast with
tangent which is increasing). Period $= \pi$. Formalize EK 3.11.A.4 and 3.11.A.5.

\columnbreak

\textbf{Step 4 --- Summary and Connections} (8 min)

Display a 3-column summary table: function, domain restriction, asymptote
equation, range, period. Students copy into guided notes.
Key comparison: sec and csc have range $(-\infty,-1]\cup[1,\infty)$; cot has
range $(-\infty,\infty)$. All three have period $2\pi$ except cot ($\pi$).
Close: \textit{``These three functions appear in calculus (especially
derivatives of trig functions) and in right-triangle trigonometry as
abbreviated notation.''} Preview homework.

\end{multicols}

\smallskip
\textbf{Pacing note:} Steps 1--2 and guided notes (Parts 1--2) in Period 1;
Steps 3--4, group activity, and exit ticket in Period 2.
\end{skillbox}

\vspace{0.08in}

%----------------------------------------------------------------------
% Active Monitoring
%----------------------------------------------------------------------
\begin{skillbox}[Active Monitoring]{redbox}
\begin{itemize}[leftmargin=1.4em, itemsep=3pt]
  \item \textbf{Confusing which denominator}: Students may mix up $\sec$ and $\csc$.
    Prompt: \textit{``Secant goes with co\underline{s}ine? No --- secant goes with
    cosine (both start with ``co''? Actually sec goes with cos, csc goes with sin.
    Mnemonic: `cosecant is the reciprocal of the shorter name, sine.'\,''}
  \item \textbf{Asymptote location}: Students may write $\theta = 0$ as an asymptote
    of sec. Correct by asking: \textit{``What is $\cos(0)$? Is it zero or one?
    So sec(0) is $1/1 = 1$, not undefined.''}
  \item \textbf{Range confusion}: Students expect range to be an interval.
    Reinforce: \textit{``$|sec\theta| \geq 1$ always. The output is never $1/2$,
    for example, because that would require $\cos\theta = 2$, which is impossible.''}
  \item \textbf{Cot increasing vs.\ decreasing}: Some students assume cot is
    increasing like tan. Show the table of cot values: as $\theta$ increases from
    $0^+$ to $\pi^-$, cot goes from $+\infty$ to $-\infty$ --- decreasing.
\end{itemize}
\end{skillbox}

\vspace{0.08in}

%----------------------------------------------------------------------
% Group Work & Differentiation
%----------------------------------------------------------------------
\begin{skillbox}[Group Work \& Differentiation]{redbox}
Students work in groups of 3--4 on the tiered activity. Teacher assigns tiers
or students self-select.

\begin{itemize}[leftmargin=1.4em, itemsep=4pt]
  \item \textbf{Tier R (Reinforce)}: Given a complete table of $\sin\theta$,
    $\cos\theta$ values at 12 key angles, students compute sec, csc, and cot
    at each angle, marking ``undefined'' where appropriate. Identify which
    entries are undefined and state the corresponding asymptote.
  \item \textbf{Tier A (Apply)}: Without scaffolding, students determine the
    full domain of sec, csc, and cot; list all vertical asymptotes in $[0,2\pi]$;
    evaluate each function at three given angles; and solve a contextual problem
    comparing the reciprocal value to the original.
  \item \textbf{Tier E (Extend)}: Students compare the graph behavior of
    $\sec\theta$ vs.\ $\cos\theta$ to explain why the arcs of sec ``open away''
    from the $\theta$-axis with vertices at the max/min of cos. Students also
    verify algebraically that $\cot\theta = \cos\theta/\sin\theta$ using
    definitions, and investigate what happens to $\cot\theta$ near its asymptotes.
\end{itemize}
\end{skillbox}

\vspace{0.08in}

%----------------------------------------------------------------------
% Individual Work & Assessment
%----------------------------------------------------------------------
\begin{skillbox}[Individual Work \& Assessment]{redbox}
Exit ticket administered independently (final 5--7 minutes of Period 2).

\begin{enumerate}[leftmargin=1.4em, itemsep=4pt]
  \item Given $\cos\theta = -1/2$, students find $\sec\theta$.
  \item Students state which $\theta$ values are excluded from the domain of
    $\csc\theta$ and explain why.
  \item Multiple-choice: identify the period of $\cot\theta$.
\end{enumerate}

\textbf{Data use:} Sort exit tickets for denominator confusion (sec vs.\ csc),
domain/asymptote errors, and range misconceptions. Use to open Lesson 3.12.
\end{skillbox}

\vspace{0.08in}

%----------------------------------------------------------------------
% Reinforcement & Extension
%----------------------------------------------------------------------
\begin{skillbox}[Reinforcement \& Extension]{goldbox}
\textbf{Homework (Lesson 7.4):} 5--6 problems covering: evaluating sec, csc,
and cot at specific unit-circle angles; identifying domain restrictions and
asymptote locations; one problem connecting the graph of $\cos\theta$ to the
behavior of $\sec\theta$ near the zeros of cosine; and 2 AP-style MC items.

\textbf{Preview of Lesson 3.12 --- Sinusoidal Function Contexts:}
Ask students: \textit{``We now have six trigonometric functions. In Lesson 3.12
we'll apply sinusoidal models to real-world contexts --- tides, temperature,
and periodic motion. As you do your homework tonight, notice how sec and csc
values become very large near their asymptotes --- that extreme behavior
appears in physics models involving resonance and oscillation.''}
\end{skillbox}

\end{document}
